% Exemplo: segunda aula do curso (subfiles), com colunas reveladas aos poucos, diagrama TikZ,
% código com trechos destacados e tabela com linhas alternadas. O preâmbulo vem de ../curso.tex.
% Compilar só esta aula: latexmk -pdf aula-02.tex   Testado com: beamer, subfiles, TeX Live 2024+
\documentclass[../curso.tex]{subfiles}

\begin{document}

\part{Como o \LaTeX\ funciona}
\begin{refsection}

\section{Compilação}

%%BEGIN:colunasreveladas
\tikzset{etapa/.style={draw=primaria, fill=primaria!10, rounded corners,
                       minimum width=2.6cm, minimum height=.8cm}}
\begin{frame}{Do texto ao PDF}
  \begin{columns}
    \begin{column}{.5\textwidth}
      \begin{itemize}[<+->]
        \item O autor escreve o arquivo \texttt{.tex}
        \item O motor (pdf\LaTeX) compõe as páginas
        \item O resultado é um PDF
      \end{itemize}
    \end{column}
    \begin{column}{.45\textwidth}
      \onslide<4->{%
      \begin{tikzpicture}[node distance=6mm]
        \node[etapa] (tex) {arquivo .tex};
        \node[etapa, below=of tex] (motor) {pdf\LaTeX};
        \node[etapa, below=of motor] (pdf) {arquivo .pdf};
        \draw[->, thick] (tex) -- (motor);
        \draw[->, thick] (motor) -- (pdf);
      \end{tikzpicture}}
    \end{column}
  \end{columns}
\end{frame}
%%END:colunasreveladas

\section{Estrutura do documento}

%%BEGIN:codigoslide
\begin{frame}[fragile]{Um documento mínimo}
\begin{lstlisting}[language={[LaTeX]TeX}]
\documentclass{article}         (*\alert{\% classe}*)
\usepackage[brazilian]{babel}   (*\alert{\% preâmbulo}*)
\begin{document}
Olá, mundo!                     (*\alert{\% corpo}*)
\end{document}
\end{lstlisting}
  O \texttt{[fragile]} é obrigatório em slides com código.
\end{frame}
%%END:codigoslide

%%BEGIN:tabelaslide
\begin{frame}{Tamanhos de fonte}
  \centering
  \rowcolors{2}{primaria!8}{white}     % linhas alternadas, a partir da segunda
  \begin{tabular}{ll}
    \toprule
    \textbf{Comando} & \textbf{Uso} \\
    \midrule
    \texttt{\textbackslash small}  & notas e legendas \\
    \texttt{\textbackslash normalsize} & texto \\
    \texttt{\textbackslash large}  & títulos curtos \\
    \texttt{\textbackslash Large}  & destaques \\
    \bottomrule
  \end{tabular}
\end{frame}
%%END:tabelaslide

\end{refsection}
\end{document}
